\section{讨论}

\subsection{有限通信下的稳定性来源}

实验结果表明，有限通信下的稳定性不仅依赖是否使用图结构，还依赖图结构是否包含任务相关的物理边语义。普通 GAT-MAPPO 能在部分半径下改善 MAPPO，但在 radius=8/10 下仍存在成功率下降和碰撞率升高问题。EA-RG-MAPPO-S 通过相对边特征和 staged 随机半径微调，使策略在多个通信半径下保持更稳定表现。

从机制上看，边特征提供了两类信息。第一类是物理几何关系，例如距离、方位和相对速度，它们直接影响追击包围与安全间隔。第二类是通信拓扑关系，例如通信可达标记和相对通信半径距离，它们反映当前信息共享边界。staged 随机半径微调进一步迫使策略在不同邻接结构下保持行为一致性，因此比只在固定半径训练更适合通信条件变化的场景。

\subsection{向 LAG/6DOF 系统扩展的边界}

本文方法的一个重要特点是表示层可扩展。当前实验使用二维追逃环境验证算法贡献，但角色图和边特征并不限定于二维状态。后续迁移到 LAG/JSBSim 或 6DOF 空战平台时，可以将节点特征扩展为三维位置、姿态、速度、过载、雷达观测和武器状态，将角色扩展为有人机、无人机、导弹、目标和僚机等类型。边特征也可以扩展为视距、雷达可探测性、导弹不可逃逸区、通信链路质量和编队任务关系。

需要强调的是，这种扩展不是直接复用全部训练环境。可复用的是角色图编码器、边特征设计思想、通信掩码机制和 MAPPO 训练接口；需要重新适配的是 6DOF 动力学、动作空间、奖励函数、传感器模型、武器交战逻辑和仿真步进接口。因此，后续工作可以沿用本文的核心算法结构，但仍需要进行平台级工程迁移和新的实验验证。当前结果不能被写成完整 6DOF 空战验证。

\subsection{局限性}

本文仍存在局限。首先，当前实验环境是二维简化追逃，不等同于完整空战系统。其次，目标意图辅助分支尚未形成可靠 balanced accuracy，因此不作为本文主贡献。最后，当前主结果使用 3 个随机种子，虽然 300 回合每种子复评已经提高可信度，但若面向更高要求期刊，仍可进一步扩展到 5 个种子或加入 LAG 小规模迁移验证。
